% Format dasar dokumen ilmiah ITB.
\usepackage{mathptmx}
\usepackage{graphicx}
\usepackage{titlesec}
\usepackage{etoolbox}

% Bab menggunakan angka Romawi dan judul huruf kapital seperti templat asal.
\renewcommand{\chaptername}{BAB}
\renewcommand{\thechapter}{\Roman{chapter}}
\titleformat{\chapter}[display]
  {\Large\centering\bfseries}
  {\chaptername\ \thechapter}{0pt}
  {\Large\bfseries\MakeUppercase}
\titleformat*{\section}{\bfseries\Large}

% Nomor gambar dan tabel mengikuti nomor subbab.
\counterwithin{figure}{section}
\counterwithin{table}{section}
\setcounter{secnumdepth}{3}

% Nama elemen dokumen dalam Bahasa Indonesia.
\renewcommand{\contentsname}{DAFTAR ISI}
\renewcommand{\listfigurename}{DAFTAR GAMBAR}
\renewcommand{\listtablename}{DAFTAR TABEL}
\renewcommand{\bibname}{DAFTAR PUSTAKA}

% Gunakan label khusus setelah Quarto mengaktifkan mode lampiran.
\apptocmd{\appendix}{\renewcommand{\chaptername}{LAMPIRAN}}{}{}

% Quarto/Pandoc masih menulis judul part lampiran PDF sebagai "Appendices".
% Normalisasi literal tersebut tanpa mengubah entri daftar isi lainnya.
\let\itbaddcontentsline\addcontentsline
\renewcommand{\addcontentsline}[3]{%
  \ifstrequal{#3}{Appendices}
    {\itbaddcontentsline{#1}{#2}{Lampiran}}
    {\itbaddcontentsline{#1}{#2}{#3}}%
}

% Sediakan ruang untuk nomor seperti III.3.1 pada daftar gambar dan tabel.
\makeatletter
\renewcommand*\l@figure{\@dottedtocline{1}{1.5em}{4em}}
\renewcommand*\l@table{\@dottedtocline{1}{1.5em}{4em}}
\makeatother

% Mencegah baris terlalu melebar pada naskah dengan banyak istilah teknis.
\setlength{\emergencystretch}{3em}
